\chapter{แผนการดำเนินงานในช่วงที่เหลือ}

\section{กรอบเวลาและเป้าหมายที่ต้องปิด}
นับจากวันที่จัดทำรายงาน (4 สิงหาคม 2569) ถึงวันสิ้นปีงบประมาณ (30 กันยายน 2569)
เหลือเวลาประมาณ 55~วัน โดยมีภาระที่ต้องปิดให้ได้ ดังนี้

\begin{table}[H]
\centering
\caption{ภาระงานที่ต้องปิดภายในสิ้นปีงบประมาณ 2569}
\small
\begin{tabular}{L{7.6cm}rr}
\toprule
\rowcolor{headerblue}
\textbf{รายการ} & \textbf{ปัจจุบัน} & \textbf{เป้าหมาย}\\
\midrule
เบิกจ่ายงบประมาณ (บาท) & 4,460,332 & 7,187,925\\
\rowcolor{rowblue}
คิดเป็นร้อยละของงบจัดสรร & 56\% & 90\%\\
วงเงินที่ต้องเบิกจ่ายเพิ่ม (บาท) & \multicolumn{2}{r}{2,727,593}\\
\rowcolor{rowblue}
รายงานผลรายกิจกรรมเข้าระบบ & 0 & 268\\
โครงการที่ต้องปิดครบ & 59 อนุมัติแล้ว & 61\\
\bottomrule
\end{tabular}
\end{table}

\section{แผนงานรายช่วงเวลา}

\begin{table}[H]
\centering
\caption{แผนการดำเนินงานสิงหาคม -- กันยายน 2569}
\small
\begin{tabular}{L{2.3cm}L{8.1cm}L{3.0cm}}
\toprule
\rowcolor{headerblue}
\textbf{ช่วงเวลา} & \textbf{กิจกรรม} & \textbf{หมุดหมาย}\\
\midrule
สิงหาคม 2569 & เร่งรัดการเบิกจ่ายและปิดใบสั่งซื้อคงค้าง ; ลงพื้นที่ติดตามผลและเก็บหลักฐานตัวชี้วัด ; สนับสนุนกระบวนการพัสดุให้หน่วยงานที่ยังไม่เบิกจ่าย & เบิกจ่ายสะสมถึงร้อยละ~75\\
\rowcolor{rowblue}
กันยายน 2569 (ต้น--กลาง) & บันทึกรายงานผลรายกิจกรรมเข้าระบบให้ครบทุกโครงการ ; ยื่นจดทะเบียนทรัพย์สินทางปัญญา ; ลงนามบันทึกความร่วมมือกับหน่วยงานภายนอก & รายงานผลเข้าระบบครบ 268~กิจกรรม\\
กันยายน 2569 (ปลาย) & ตรวจรับผลตัวชี้วัดจากหลักฐานจริง ; จัดทำรายงานผลการดำเนินงานฉบับสมบูรณ์ ; ถอดบทเรียนและจัดทำข้อเสนอปีงบประมาณ 2570 & ปิดปีงบประมาณ 30 กันยายน 2569\\
\bottomrule
\end{tabular}
\end{table}

\section{มาตรการเร่งรัด 4 ด้าน}

\subsection{มาตรการที่ 1 --- เร่งรัดการเบิกจ่าย}
ติดตามรายโครงการที่เบิกจ่ายต่ำกว่าร้อยละ~30 เป็นรายสัปดาห์ พร้อมจัดทีมสนับสนุน
ช่วยจัดทำเอกสารจัดซื้อจัดจ้างให้หน่วยงานที่ยังไม่มีการเบิกจ่าย โดยเฉพาะ มทร.ล้านนา ตาก
(3~โครงการ 250,000~บาท) และคณะบริหารธุรกิจและศิลปศาสตร์ (1~โครงการ 40,000~บาท)

\subsection{มาตรการที่ 2 --- เก็บผลงานเข้าระบบ}
กำหนดให้ผู้รับผิดชอบโครงการรายงานผลรายกิจกรรมพร้อมหลักฐานภาพถ่ายและเอกสาร
ในระบบติดตามโครงการ แล้วให้กลุ่มแผนงานตรวจรับและยืนยันผลตัวชี้วัดจากหลักฐานจริง
มาตรการนี้เป็นเงื่อนไขจำเป็นของการจัดทำรายงานฉบับสมบูรณ์

\subsection{มาตรการที่ 3 --- ปิดช่องว่างตัวชี้วัด}
เร่งยื่นจดทะเบียนทรัพย์สินทางปัญญาจากผลงานที่พร้อม (ตัวชี้วัดที่ 17)
และจัดทำบันทึกความร่วมมือกับหน่วยงานภายนอกและสถานประกอบการในพื้นที่
(ตัวชี้วัดที่ 10 และที่ 38) โดยใช้ฐานความสัมพันธ์ที่มีอยู่แล้วจากการทำงานในพื้นที่

\subsection{มาตรการที่ 4 --- เตรียมต่อยอดปีงบประมาณ 2570}
ถอดบทเรียนรายพื้นที่และรายหน่วยงาน แล้วออกแบบตัวชี้วัดของแผนงานให้สอดคล้องกับ
ลักษณะงานพัฒนาชุมชนตั้งแต่ต้นปีงบประมาณ พร้อมกำหนดปฏิทินการเบิกจ่ายรายไตรมาส
เพื่อไม่ให้เกิดการกระจุกตัวของการเบิกจ่ายในช่วงปลายปีอีก

\section{การตอบโจทย์ในภาพกว้าง}
ผลการดำเนินงานของแผนงานตอบโจทย์ผู้มีส่วนได้ส่วนเสียใน 4 ระดับ ดังนี้

\begin{table}[H]
\centering
\caption{การตอบโจทย์ผู้มีส่วนได้ส่วนเสีย 4 ระดับ}
\small
\begin{tabular}{L{3.0cm}L{10.7cm}}
\toprule
\rowcolor{headerblue}
\textbf{ระดับ} & \textbf{การตอบโจทย์}\\
\midrule
มหาวิทยาลัย & ส่งมอบตัวชี้วัดตามยุทธศาสตร์ความเป็นเลิศ 7~ตัว และทำเป้าภารกิจโดยตรง (ตัวชี้วัดที่ 39) ได้เกินร้อยละ~100\\
\rowcolor{rowblue}
หน่วยงานในพื้นที่ & ศูนย์และสถานีโครงการหลวง โครงการตามพระราชดำริ โรงเรียน วัด และวิสาหกิจชุมชน ได้รับเทคโนโลยี องค์ความรู้ และกำลังคนที่ใช้งานได้จริง\\
ประเทศ & ลดความเหลื่อมล้ำบนพื้นที่สูง สร้างอาชีพ และเพิ่มขีดความสามารถของเกษตรกรและชุมชน\\
\rowcolor{rowblue}
ระดับโลก & เชื่อมโยงกับเป้าหมายการพัฒนาที่ยั่งยืนขององค์การสหประชาชาติ\\
\bottomrule
\end{tabular}
\end{table}

ในภาพรวมระดับแผนงาน การดำเนินงานเชื่อมโยงกับเป้าหมายการพัฒนาที่ยั่งยืนหลัก ได้แก่
SDG 1 ขจัดความยากจน $\cdot$ SDG 2 ยุติความหิวโหย $\cdot$ SDG 4 การศึกษาที่มีคุณภาพ $\cdot$
SDG 8 งานที่มีคุณค่าและการเติบโตทางเศรษฐกิจ $\cdot$ SDG 9 อุตสาหกรรม นวัตกรรม และโครงสร้างพื้นฐาน $\cdot$
SDG 10 ลดความเหลื่อมล้ำ $\cdot$ SDG 13 การรับมือการเปลี่ยนแปลงสภาพภูมิอากาศ $\cdot$
SDG 17 ความร่วมมือเพื่อการพัฒนาที่ยั่งยืน
และเป้าหมายสนับสนุน ได้แก่ SDG 7 พลังงานสะอาด $\cdot$ SDG 11 ชุมชนยั่งยืน $\cdot$
SDG 12 การผลิตและบริโภคที่ยั่งยืน $\cdot$ SDG 15 ระบบนิเวศบนบก

\begin{infobox}
การเชื่อมโยงเป้าหมายการพัฒนาที่ยั่งยืนข้างต้นเป็นการวิเคราะห์\textbf{ระดับแผนงาน}
จากวัตถุประสงค์ของโครงการหลัก ยังมิได้ผูกรายโครงการไว้ในฐานข้อมูล
การผูก SDG รายโครงการเป็นงานที่ควรดำเนินการในปีงบประมาณ 2570
\end{infobox}
